\documentclass[11pt,a4paper]{article}
\usepackage[margin=2.3cm]{geometry}
\usepackage{booktabs}
\usepackage{graphicx}
\usepackage{amsmath}
\usepackage{caption}
\usepackage{pgfplots}
\pgfplotsset{compat=1.17}

\pgfplotscreateplotcyclelist{tenseries}{
  color=red,mark=*\\
  color=blue,mark=square*\\
  color=green!50!black,mark=triangle*\\
  color=orange,mark=diamond*\\
  color=violet,mark=pentagon*\\
  color=cyan!70!black,mark=o\\
  color=magenta,mark=square\\
  color=brown,mark=triangle\\
  color=gray,mark=diamond\\
  color=black,mark=x\\
}
\pgfplotsset{
  report plot/.style={
    width=\linewidth, height=7.5cm, grid=major, cycle list name=tenseries,
    legend columns=5, legend style={at={(0.5,-0.3)}, anchor=north, font=\scriptsize}
  }
}
\newcommand{\fit}[1]{\resizebox{\ifdim\width>\linewidth\linewidth\else\width\fi}{!}{#1}}
\title{\textbf{Lab Report} \\ Operating Systems Laboratory \\ Assignment 5: Parallel Applications}
\author{Anshuman Rahul (CS24BT060) \\ Sidharth B (CS24BT061)}
\date{}

% ---------------------------------------------------------------------
% Additions for this report
% ---------------------------------------------------------------------
\usepgfplotslibrary{groupplots}     % grids of plots (Figures 3 and 4)
\usepackage{pgfplotstable}          % computed columns (relative times in Figures 3 and 4)
\usepackage[section]{placeins}      % tables/figures stay within their own section
\usepackage[hidelinks]{hyperref}    % clickable cross-references

% Run configuration (config.hpp).
\newcommand{\cfgIterations}{1000}
\newcommand{\cfgK}{8}
\newcommand{\cfgRing}{16}
\newcommand{\cfgPasses}{100}
\newcommand{\cfgResults}{10}

\pgfplotscreateplotcyclelist{programs}{
  color=red,mark=*\\
  color=blue,mark=square*\\
  color=green!50!black,mark=triangle*\\
}
\pgfplotscreateplotcyclelist{rings}{
  color=red,mark=*\\
  color=green!50!black,mark=triangle*\\
}
\pgfplotsset{
  sweep plot/.style={
    width=0.47\linewidth, height=5cm, grid=major,
    xmode=log, log basis x=2, log ticks with fixed point,
    tick label style={font=\scriptsize}, label style={font=\small},
    title style={font=\small}
  }
}
\newcommand{\code}[1]{\texttt{#1}}

% ---------------------------------------------------------------------
% Plot data (seconds). Columns: part2_1 = threads, part2_2 = pipes, part2_3 = shared memory.
% ---------------------------------------------------------------------

% Speedup over serial at k = 8, N = 16; bound = T_serial / T_S1 (serial run).
\pgfplotstableread{
image part2_1 part2_2 part2_3 bound
1 0.9798 0.9592 0.9801 1.0860
2 0.9603 0.9317 0.9464 1.0790
4 0.9491 0.9234 0.9434 1.0760
}\speedupVOne
\pgfplotstableread{
image part2_1 part2_2 part2_3 bound
1 1.0768 1.0676 1.0659 1.0720
2 1.0766 1.0675 1.0724 1.0732
4 1.0748 1.0650 1.0623 1.0746
}\speedupVTwo

% Pipeline time against k (N = 16).
\pgfplotstableread{
k part2_1 part2_2 part2_3
1 3.05 3.27 3.31
2 3.03 3.22 3.41
4 2.97 3.07 3.10
8 2.98 3.05 2.99
16 2.99 3.05 3.08
32 3.01 3.13 2.96
64 2.99 4.32 2.98
128 2.98 5.45 2.95
}\sweepBimgOneVOne
\pgfplotstableread{
k part2_1 part2_2 part2_3
1 2.62 2.66 2.72
2 2.67 2.65 2.65
4 2.64 2.65 2.66
8 2.63 2.64 2.70
16 2.61 2.65 2.68
32 2.64 2.64 2.63
64 2.65 2.76 2.68
128 2.64 2.71 2.65
}\sweepBimgOneVTwo
\pgfplotstableread{
k part2_1 part2_2 part2_3
1 11.34 11.93 11.81
2 11.30 11.55 11.72
4 11.32 11.60 11.45
8 11.29 11.57 11.47
16 11.28 11.62 11.40
32 11.23 16.34 11.37
64 11.25 15.39 11.39
128 11.30 13.73 11.31
}\sweepBimgTwoVOne
\pgfplotstableread{
k part2_1 part2_2 part2_3
1 9.98 10.09 10.08
2 9.96 10.10 10.05
4 10.01 10.03 9.99
8 9.93 10.03 9.99
16 9.96 10.04 10.07
32 10.00 10.11 10.05
64 10.01 10.03 10.03
128 9.95 10.06 10.02
}\sweepBimgTwoVTwo
\pgfplotstableread{
k part2_1 part2_2 part2_3
1 42.77 43.70 43.95
2 42.55 43.62 43.40
4 42.60 43.28 43.17
8 42.61 43.47 43.01
16 42.55 43.77 42.87
32 42.61 57.67 42.66
64 42.69 60.43 42.77
128 42.59 56.17 42.78
}\sweepBimgFourVOne
\pgfplotstableread{
k part2_1 part2_2 part2_3
1 37.85 37.78 37.78
2 37.59 37.97 37.72
4 37.70 37.84 37.92
8 37.66 37.94 37.80
16 37.64 37.96 37.91
32 37.63 37.90 37.75
64 37.77 37.84 37.74
128 37.75 37.87 37.83
}\sweepBimgFourVTwo

% Pipeline time against N (k = 8).
\pgfplotstableread{
N part2_1 part2_3
1 3.78 12.00
2 2.99 3.09
4 2.94 3.04
8 2.97 3.04
16 2.98 3.00
32 2.98 3.02
64 2.98 3.00
}\sweepCimgOneVOne
\pgfplotstableread{
N part2_1 part2_3
1 2.65 3.01
2 2.64 2.71
4 2.65 2.69
8 2.62 2.64
16 2.66 2.70
32 2.62 2.65
64 2.66 2.65
}\sweepCimgOneVTwo
\pgfplotstableread{
N part2_1 part2_3
1 14.30 42.95
2 11.39 13.51
4 11.19 11.46
8 11.29 11.36
16 11.24 11.39
32 11.28 11.41
64 11.28 11.45
}\sweepCimgTwoVOne
\pgfplotstableread{
N part2_1 part2_3
1 10.02 10.28
2 10.05 10.03
4 10.00 10.05
8 10.03 10.00
16 10.00 10.01
32 10.01 10.09
64 9.98 10.01
}\sweepCimgTwoVTwo
\pgfplotstableread{
N part2_1 part2_3
1 54.39 129.86
2 42.36 43.35
4 42.37 42.81
8 42.33 43.02
16 42.69 43.00
32 42.61 43.06
64 42.50 42.95
}\sweepCimgFourVOne
\pgfplotstableread{
N part2_1 part2_3
1 37.79 38.92
2 37.61 37.93
4 37.93 37.89
8 37.62 38.06
16 37.95 37.99
32 37.83 37.79
64 37.74 37.73
}\sweepCimgFourVTwo

\begin{document}
\maketitle

% =====================================================================
\section{Introduction}
The application is the image sharpener from Lab~1. It runs three stages on every pixel and
colour channel of a PPM image:
\begin{itemize}
\item \textbf{S1, smooth:} the mean of the $3\times3$ neighbourhood (edge pixels are replicated at
  the borders);
\item \textbf{S2, details:} $d = \max(0,\ \text{original} - \text{smoothed})$;
\item \textbf{S3, sharpen:} $\text{output} = \min(255,\ \text{original} + d)$.
\end{itemize}
S2 needs S1's result and S3 needs S2's, so the stages can run in parallel only as a
\emph{pipeline}: while S3 works on one part of the image, S2 works on the next and S1 on the one
after that. This report covers a serial version (Part~I); three pipelined versions on one
machine, using threads with atomics, processes with pipes, and processes with shared memory and
semaphores (Part~II); and a version split over two machines connected by TCP (Part~III). Every
pipelined version is built in two variants (Section~\ref{sec:variants}). For each version we
check that the data arrives exactly as sent, and we measure runtime and speedup.

% =====================================================================
\section{Experimental setup}

\subsection{The two variants}
\label{sec:variants}
\begin{itemize}
\item \textbf{Variant 1:} the whole computation (S1, S2, S3 over the image) is repeated
  \cfgIterations{} times, and every iteration streams through the pipeline. The output of the last
  iteration is written to the file.
\item \textbf{Variant 2:} every stage runs \cfgIterations{} passes in blocks of \cfgPasses. Only
  one result per block is handed to the next stage, so each link carries \cfgResults{} images
  instead of \cfgIterations. The other passes are ordinary computations whose output stays inside
  the stage.
\end{itemize}
Every pass reads unchanged input, so all iterations and passes produce the same image, and the
final output equals one run of the serial sharpener.

\subsection{Run configuration and method}
\begin{center}
\begin{tabular}{lr}
\toprule
Parameter & Value \\
\midrule
Iterations (Variant 1) / passes per stage (Variant 2) & \cfgIterations \\
Passes per result (Variant 2) & \cfgPasses \\
Results sent per link (Variant 2) & \cfgResults \\
Default strip height $k$ (rows) & \cfgK \\
Default ring capacity $N$ (strips) & \cfgRing \\
Runs per configuration & 1 \\
\bottomrule
\end{tabular}
\end{center}
All programs are compiled with \code{g++ -O2} and print their own measurements. All times come
from \code{std::chrono::steady\_clock}:
\begin{itemize}
\item \textbf{Serial time} is the time spent computing all passes of S1, S2 and S3 (excluding
  reading and writing the image).
\item \textbf{Pipeline time} runs from starting the stages until all of them have finished,
  including S3 writing the output file (a few milliseconds at most). Reading the input image comes
  before and is not included.
\item \textbf{Speedup} is serial time divided by pipeline time, for the same image and variant.
\end{itemize}
Each stage also keeps its own counters (busy, verify, wait in, wait out; see
Section~\ref{sec:stages}), and each end of every link counts what it sent or received
(Section~\ref{sec:verification}). Each configuration was run once. As a rough guide to
run-to-run variation, in Variant~2 the strip height should not matter at all, yet the threads
version's times for image~4 vary by 0.7\% across the values of $k$ (Table~\ref{tab:k-v2}).
Differences of about 1\% are therefore not significant.

\subsection{Test images}
All measurements use the same three images (Table~\ref{tab:images}, Figure~\ref{fig:images}).
\begin{table}[htbp]
\centering
\caption{Test images. Strips per image are for the default $k=\cfgK$.}
\label{tab:images}
\begin{tabular}{rrrrrr}
\toprule
Image & Width & Height & Pixels & File size (MB) & Strips per image \\
\midrule
1 & 450 & 180 & 81{,}000 & 0.24 & 23 \\
2 & 786 & 393 & 308{,}898 & 0.93 & 50 \\
4 & 1280 & 910 & 1{,}164{,}800 & 3.49 & 114 \\
\bottomrule
\end{tabular}
\end{table}

\begin{figure}[p]
\centering
\newcommand{\imagepair}[1]{%
  \begin{minipage}[t]{0.48\linewidth}\centering
    \includegraphics[width=\linewidth]{figures/img_#1_in.png}\\[2pt]
    \small Image #1, input
  \end{minipage}\hfill
  \begin{minipage}[t]{0.48\linewidth}\centering
    \includegraphics[width=\linewidth]{figures/img_#1_out.png}\\[2pt]
    \small Image #1, sharpened
  \end{minipage}}
\imagepair{1}\par\bigskip
\imagepair{2}\par\bigskip
\imagepair{4}
\caption{The test images and the sharpened outputs. Every program produced exactly these outputs
(byte-identical to the serial reference).}
\label{fig:images}
\end{figure}

% =====================================================================
\section{Design and implementation}

\subsection{Common structure}
\paragraph{Strips.} The unit of work and of communication is a \emph{strip} of $k$ consecutive
image rows ($k$ is a command-line argument, default \cfgK). The last strip of an image has
$h \bmod k$ rows when $k$ does not divide the height $h$. All three kernels work on row ranges:
\code{S1\_smooth(orig, r0, r1, out)} computes rows $[r_0, r_1)$, reading one row above and below
from the original image, and \code{S2\_details}/\code{S3\_sharpen} work pointwise on $n$ bytes.
A stage can therefore compute a strip as soon as it has received it, and the pipeline fills
after one strip instead of one full image. Images are stored as one flat array of
$h \times w \times 3$ bytes, so a strip is one contiguous block of memory.

\paragraph{Strip header.} Every strip travels with a 24-byte header: \code{iteration} (the
iteration in Variant~1, the result number in Variant~2), \code{r0} (first row), \code{nrows},
\code{flags} (reserved, 0) and a 64-bit \code{checksum} (Section~\ref{sec:verification}).

\paragraph{Channel interface.} All three transports offer the same four operations, so the stage
code is the same apart from the channel type:
\begin{itemize}
\item producer: \code{acquire()} returns a slot (header and pixel area) to compute into;
  \code{publish()} hands it to the consumer;
\item consumer: \code{receive()} returns the next filled slot; \code{release()} gives it back once
  it is no longer needed.
\end{itemize}
S1 computes straight into the acquired slot. S2 receives a strip, computes the details into a
slot of its output channel, publishes that, and releases its input. S3 computes into the final
image.

\subsection{Part I: serial version}
Part~I runs S1, S2 and S3 once over the whole image and writes the result; its output is the
reference for all other parts. For the speedup measurements, a separate serial baseline does
exactly the work of the pipelined versions (\cfgIterations{} passes of each stage, in both
variants, with the same kernels) in one thread, and times each stage separately.

\subsection{Part II.1: threads and atomics}
The three stages are three threads of one process, connected by two single-producer,
single-consumer ring buffers of $N$ slots (default \cfgRing). Each ring holds two 64-bit counters
that only ever grow: \code{head} (slots published, written only by the producer) and \code{tail}
(slots released, written only by the consumer). Slot $i$ is at index $i \bmod N$; the ring is
full when $\text{head}-\text{tail} = N$ and empty when $\text{head} = \text{tail}$.
\begin{itemize}
\item \code{publish()} stores \code{head + 1} with \emph{release} ordering, and \code{receive()}
  loads \code{head} with \emph{acquire} ordering. All writes to the slot's header and pixels
  therefore happen-before the consumer reads them.
\item In the other direction, \code{release()} stores \code{tail} with release ordering and
  \code{acquire()} loads it with acquire ordering. The producer therefore never overwrites a slot
  the consumer is still reading.
\end{itemize}
No locks are used. \code{head} and \code{tail} are on separate cache lines, so the two threads
do not slow each other down through false sharing. A thread that has to wait spins with a
\code{pause} instruction for up to 1024 rounds and then calls \code{sched\_yield()}.

\subsection{Part II.2: processes and pipes}
The three stages are three processes created with \code{fork()}, connected by two pipes created
before the fork. Each child closes the pipe ends it does not use, and the parent closes all of
them. This matters for failures: a reader only sees end-of-file, and a writer only gets
\code{EPIPE}, once \emph{every} copy of the other end is closed. A pipe is a byte stream, so a
strip is sent as its header followed by $\text{nrows}\times w\times3$ pixel bytes. The receiver
first reads the header, checks that \code{nrows} is at most $k$, and then reads exactly that many
pixel bytes. \code{readn}/\code{writen} loop until all bytes are transferred, because a pipe may
return fewer bytes than requested, and they retry after \code{EINTR}. Unlike the rings, every
strip is copied twice: from the producer into the kernel and from the kernel into the consumer.
The parent waits for the three children and reports a failure if any of them did not exit
normally.

\subsection{Part II.3: processes, shared memory and semaphores}
The stages are again three processes, but they communicate through two rings in POSIX shared
memory. Each ring is one region created with \code{shm\_open}, sized with \code{ftruncate} and
mapped with \code{mmap(MAP\_SHARED)} by the parent \emph{before} forking, so all children
inherit the mapping. The name is unlinked immediately after mapping, so nothing is left behind in
\code{/dev/shm}, even if a process crashes. The region holds two unnamed process-shared
semaphores, $N$ headers and $N$ pixel slots:
\begin{itemize}
\item \code{empty} (initially $N$): \code{acquire()} waits on it, and \code{release()} posts it;
\item \code{full} (initially 0): \code{receive()} waits on it, and \code{publish()} posts it.
\end{itemize}
The semaphores provide both the waiting and the memory ordering. Each process keeps its own
slot index. The parent waits for whichever stage finishes first; if a stage fails, the parent
kills the remaining ones, which would otherwise wait on a semaphore forever.

\subsection{Variant 2}
Every stage loops over \cfgResults{} results of \cfgPasses{} passes each. Only pass~0 of each
block communicates: S1 computes into channel slots and publishes them, S2 and S3 receive. On the
other passes, S1 writes into a private scratch buffer. S2 and S3 keep the strips received on
pass~0 in a private buffer and recompute from that unchanged input on the remaining passes. Every
stage thus still does \cfgIterations{} passes of work, as in Variant~1 and the serial baseline,
but each link carries only \cfgResults{} images.

\subsection{Part III: two machines}
The pipeline is split over two machines, A and B, connected by one TCP connection:
\begin{itemize}
\item \textbf{Partition 1:} A runs S1 and S2 as two threads with a ring between them (as in
  Part~II.1); B runs S3. The S2$\to$S3 link crosses the network.
\item \textbf{Partition 2:} A runs S1; B runs S2 and S3 as two threads with a ring. The
  S1$\to$S2 link crosses the network.
\end{itemize}
The socket uses the same byte-stream channel as the pipes. Header fields are converted to network
byte order, and Nagle's algorithm is disabled (\code{TCP\_NODELAY}) so strips are sent
immediately. B is started first and listens; A connects, retrying for about 10 seconds. A then
sends a fixed-size run-parameter message:
\begin{itemize}
\item a magic number, the image size, $k$, the variant, the ring capacity;
\item the iteration and pass counts compiled into A;
\item a 64-bit hash of A's input image.
\end{itemize}
B takes $k$ (and, in partition~2, $N$) from this message, and refuses to start if the size, variant, counts or image
hash differ from its own. The two machines' clocks are not synchronised, so all timing is done
on A: after writing the output file, B sends a one-byte ``done'' message, and A's end-to-end
time runs from starting its stages until that byte arrives.

% =====================================================================
\section{Verifying correctness}
\label{sec:verification}
Correctness is checked at four levels. All checks are always on, including in the measured
runs, so the reported times include their cost.

\paragraph{1. Output against the reference.} Every output image is compared byte by byte (with
\code{cmp}) against the output of the serial sharpener. This shows end to end that the stages
computed the right result and that S3 put every strip in the right place.

\paragraph{2. Per-strip header and checksum.} Just before publishing a strip, the sender
computes a 64-bit FNV-1a hash over the header fields and the pixel bytes and stores it in the
header. Every stage walks the same fixed loop over iterations and row ranges, so the receiver
knows exactly which strip must come next. For each strip, \code{verify\_strip} checks that
\code{iteration}, \code{r0} and \code{nrows} have the expected values, and recomputes and compares
the checksum. On any mismatch, the program prints which strip failed and aborts before writing
an output. These checks catch:
\begin{itemize}
\item corrupted pixel or header bytes (for example a slot overwritten too early, or a partial
  read), through a checksum mismatch;
\item lost, duplicated or reordered strips, through an unexpected position;
\item a stream that ends too early because a stage died (pipes and sockets), through
  end-of-file in the middle of a strip.
\end{itemize}

\paragraph{3. Link statistics on both ends.} Both ends of every link count the strips and
payload bytes that pass through. They also keep a running hash, \code{stream}, into which each
strip's checksum is folded in order. The sender updates its counters on publishing. The
receiver updates its counters only after \code{verify\_strip} has accepted the strip. Equal
values on both sides show that the receiver accepted exactly the strips that were sent, with the
same content and in the same order: swapping two strips changes the stream hash. The expected
counts follow from the image: $\lceil h/k\rceil$ strips per image, times \cfgIterations{}
(Variant~1) or \cfgResults{} (Variant~2) images.

\paragraph{4. Run parameters (Part~III).} The handshake described above ensures that both
machines work on the same image with the same settings. This matters because S2 and S3 on B
read B's own copy of the input image, which the per-strip checks alone would not cover.

\paragraph{Limitations.} FNV-1a is not a cryptographic hash. It detects accidental corruption,
not deliberate tampering; an accidental change goes unnoticed with probability about $2^{-64}$
per strip. The strip checks cover the \emph{transfer} only; that the stages \emph{compute}
correctly is shown by the comparison with the reference.

\paragraph{Results.} Table~\ref{tab:links} shows both ends of every link for image~4. In Variant~1 each link carries
$114\times1000 = 114{,}000$ strips and in Variant~2 $114\times10 = 1{,}140$, as expected.

\begin{table}[htbp]
\centering
\caption{Link statistics of both ends of every link (image 4, $k=\cfgK$, $N=\cfgRing$). Strips
and bytes are totals over the run.}
\label{tab:links}
\fit{%
\begin{tabular}{llllrrl}
\toprule
Variant & Program & Link & Side & Strips & Payload bytes & Stream hash \\
\midrule
1 & Threads & S1$\to$S2 & send & 114{,}000 & 3{,}494{,}400{,}000 & \texttt{29d8d98fb66b61fd} \\
 &  & S1$\to$S2 & recv & 114{,}000 & 3{,}494{,}400{,}000 & \texttt{29d8d98fb66b61fd} \\
 &  & S2$\to$S3 & send & 114{,}000 & 3{,}494{,}400{,}000 & \texttt{d178aee5e5809a0b} \\
 &  & S2$\to$S3 & recv & 114{,}000 & 3{,}494{,}400{,}000 & \texttt{d178aee5e5809a0b} \\
\addlinespace
1 & Pipes & S1$\to$S2 & send & 114{,}000 & 3{,}494{,}400{,}000 & \texttt{29d8d98fb66b61fd} \\
 &  & S1$\to$S2 & recv & 114{,}000 & 3{,}494{,}400{,}000 & \texttt{29d8d98fb66b61fd} \\
 &  & S2$\to$S3 & send & 114{,}000 & 3{,}494{,}400{,}000 & \texttt{d178aee5e5809a0b} \\
 &  & S2$\to$S3 & recv & 114{,}000 & 3{,}494{,}400{,}000 & \texttt{d178aee5e5809a0b} \\
\addlinespace
1 & Shared mem. & S1$\to$S2 & send & 114{,}000 & 3{,}494{,}400{,}000 & \texttt{29d8d98fb66b61fd} \\
 &  & S1$\to$S2 & recv & 114{,}000 & 3{,}494{,}400{,}000 & \texttt{29d8d98fb66b61fd} \\
 &  & S2$\to$S3 & send & 114{,}000 & 3{,}494{,}400{,}000 & \texttt{d178aee5e5809a0b} \\
 &  & S2$\to$S3 & recv & 114{,}000 & 3{,}494{,}400{,}000 & \texttt{d178aee5e5809a0b} \\
\addlinespace
2 & Threads & S1$\to$S2 & send & 1{,}140 & 34{,}944{,}000 & \texttt{d20f58811324790b} \\
 &  & S1$\to$S2 & recv & 1{,}140 & 34{,}944{,}000 & \texttt{d20f58811324790b} \\
 &  & S2$\to$S3 & send & 1{,}140 & 34{,}944{,}000 & \texttt{828a6fd8aeb69287} \\
 &  & S2$\to$S3 & recv & 1{,}140 & 34{,}944{,}000 & \texttt{828a6fd8aeb69287} \\
\addlinespace
2 & Pipes & S1$\to$S2 & send & 1{,}140 & 34{,}944{,}000 & \texttt{d20f58811324790b} \\
 &  & S1$\to$S2 & recv & 1{,}140 & 34{,}944{,}000 & \texttt{d20f58811324790b} \\
 &  & S2$\to$S3 & send & 1{,}140 & 34{,}944{,}000 & \texttt{828a6fd8aeb69287} \\
 &  & S2$\to$S3 & recv & 1{,}140 & 34{,}944{,}000 & \texttt{828a6fd8aeb69287} \\
\addlinespace
2 & Shared mem. & S1$\to$S2 & send & 1{,}140 & 34{,}944{,}000 & \texttt{d20f58811324790b} \\
 &  & S1$\to$S2 & recv & 1{,}140 & 34{,}944{,}000 & \texttt{d20f58811324790b} \\
 &  & S2$\to$S3 & send & 1{,}140 & 34{,}944{,}000 & \texttt{828a6fd8aeb69287} \\
 &  & S2$\to$S3 & recv & 1{,}140 & 34{,}944{,}000 & \texttt{828a6fd8aeb69287} \\
\bottomrule
\end{tabular}}
\end{table}

% =====================================================================
\section{Results: Part II}
\label{sec:results2}
All three images were run in both variants. Each table states the strip height $k$ and ring
capacity $N$ it was measured with; the defaults are $k=\cfgK$ and $N=\cfgRing$. Where a best
value is meaningful (the fastest program, or the best $k$ or $N$), it is shown in bold.

\subsection{How fast can the pipeline be?}
The serial program runs the three stages one after the other, so its time is the sum of the
three stage times:
\[
  T_\text{serial} = T_\text{S1} + T_\text{S2} + T_\text{S3}.
\]
In the pipeline, every stage is a separate thread or process, and S1 performs exactly the same
computation as in the serial program: all \cfgIterations{} passes over the image. S1 is a single
thread of execution, so however fast S2 and S3 are, the pipeline cannot finish before S1 has done
all of its work. If $T'_\text{S1}$ is the time S1 spends working in a pipeline run (computing
and checksums), then
\begin{equation}
  T_\text{pipeline} \;\ge\; T'_\text{S1}.
  \label{eq:s1}
\end{equation}
If S1 also works as fast in the pipeline as in the serial program ($T'_\text{S1} \approx
T_\text{S1}$), this gives an upper bound for the speedup:
\begin{equation}
  \text{speedup} = \frac{T_\text{serial}}{T_\text{pipeline}}
  \;\le\; \frac{T_\text{serial}}{T_\text{S1}}.
  \label{eq:bound}
\end{equation}
This bound rests on three assumptions:
\begin{enumerate}
\item every stage has a core of its own, so the stages really run at the same time (the machines
  used have more cores than there are stages);
\item each stage does the same work in both programs (true here: same kernels, same number of
  passes);
\item each stage works at the same speed in both programs. This holds only approximately: cache
  behaviour and clock speed differ slightly from run to run, so a single run can end up marginally
  above the bound.
\end{enumerate}
It is also an ideal case: it ignores everything the pipeline adds (communication, checksums, and
the roughly 1\,ms until the first strip reaches S3), so real runs lie below it.

Inequality~\eqref{eq:s1} needs no assumptions about the serial program, and both sides are measured
in the same run, so it can be checked directly. For image~4, Variant~2, with threads,
S1 spent $37.46 + 0.05 = 37.51$\,s computing and checksumming (Table~\ref{tab:stages4}), and the
pipeline took 37.51\,s: the pipeline finished as soon as S1 did.

Table~\ref{tab:serial} shows the stage times measured in the serial program. For image~4 in
Variant~1, for example, $T_\text{serial} = 37.63 + 1.32 + 1.54 = 40.49$\,s, so the bound is
$40.49 / 37.63 = 1.076$. S1 averages nine input values for every output value, while S2 and S3
only combine two, so S1 takes over 90\% of the time and the bound is only about $1.08\times$,
however fast the stages communicate.
\begin{table}[htbp]
\centering
\caption{Time spent in each stage by the serial program (seconds, all \cfgIterations{} passes),
and the resulting speedup bound.}
\label{tab:serial}
\begin{tabular}{rrrrrrr}
\toprule
Image & Variant & $T_\text{S1}$ & $T_\text{S2}$ & $T_\text{S3}$ & $T_\text{serial}$ &
Bound $T_\text{serial}/T_\text{S1}$ \\
\midrule
1 & 1 & 2.70 & 0.11 & 0.12 & 2.93 & 1.086 \\
1 & 2 & 2.67 & 0.09 & 0.10 & 2.86 & 1.072 \\
\addlinespace
2 & 1 & 10.04 & 0.37 & 0.42 & 10.83 & 1.079 \\
2 & 2 & 10.01 & 0.33 & 0.40 & 10.74 & 1.073 \\
\addlinespace
4 & 1 & 37.63 & 1.32 & 1.54 & 40.49 & 1.076 \\
4 & 2 & 37.52 & 1.26 & 1.54 & 40.32 & 1.075 \\
\bottomrule
\end{tabular}
\end{table}

\subsection{Runtime and speedup}
\begin{table}[htbp]
\centering
\caption{Runtime (seconds) and speedup over the serial version for the default $k=\cfgK$,
$N=\cfgRing$. Serial = compute time; the others = pipeline time. Best per row in bold.}
\label{tab:speedup}
\fit{%
\begin{tabular}{rrr rr rr rr}
\toprule
& & & \multicolumn{2}{c}{Threads} & \multicolumn{2}{c}{Pipes} & \multicolumn{2}{c}{Shared mem.} \\
\cmidrule(lr){4-5}\cmidrule(lr){6-7}\cmidrule(lr){8-9}
Image & Variant & Serial & Time & Speedup & Time & Speedup & Time & Speedup \\
\midrule
1 & 1 & 2.93 & \textbf{2.99} & \textbf{0.980} & 3.05 & 0.959 & \textbf{2.99} & \textbf{0.980} \\
1 & 2 & 2.86 & \textbf{2.66} & \textbf{1.077} & 2.68 & 1.068 & 2.68 & 1.066 \\
\addlinespace
2 & 1 & 10.83 & \textbf{11.28} & \textbf{0.960} & 11.62 & 0.932 & 11.44 & 0.946 \\
2 & 2 & 10.74 & \textbf{9.98} & \textbf{1.077} & 10.06 & 1.067 & 10.01 & 1.072 \\
\addlinespace
4 & 1 & 40.49 & \textbf{42.66} & \textbf{0.949} & 43.85 & 0.923 & 42.92 & 0.943 \\
4 & 2 & 40.32 & \textbf{37.51} & \textbf{1.075} & 37.86 & 1.065 & 37.96 & 1.062 \\
\bottomrule
\end{tabular}}
\end{table}

\begin{figure}[htbp]
\centering
\begin{tikzpicture}
\begin{groupplot}[
  group style={group size=2 by 1, horizontal sep=1.5cm},
  width=0.5\linewidth, height=6cm,
  ybar, /pgf/bar width=5pt,
  symbolic x coords={1,2,4}, xtick=data, enlarge x limits=0.25,
  xlabel={Image}, ymin=0, ymax=1.2, ytick={0,0.2,...,1.2},
  ymajorgrids, extra y ticks={1}, extra y tick labels={},
  extra y tick style={grid=major, major grid style={black, dashed}},
  tick label style={font=\scriptsize}, label style={font=\small},
  title style={font=\small},
  legend columns=4, legend style={font=\small, /tikz/every even column/.append style={column sep=6pt}}
]
\nextgroupplot[title={Variant 1}, ylabel={Speedup}, legend to name=legSpeedup]
\addplot[fill=red!60, /pgf/bar shift=-6pt] table[x=image, y=part2_1] {\speedupVOne};
\addplot[fill=blue!60, /pgf/bar shift=0pt] table[x=image, y=part2_2] {\speedupVOne};
\addplot[fill=green!50!black!60, /pgf/bar shift=6pt] table[x=image, y=part2_3] {\speedupVOne};
\addplot[only marks, mark=-, mark size=10pt, very thick, black, /pgf/bar shift=0pt,
         legend image code/.code={\draw[very thick] (0cm,0cm) -- (0.4cm,0cm);}]
  table[x=image, y=bound] {\speedupVOne};
\legend{Threads, Pipes, Shared mem., Bound}
\nextgroupplot[title={Variant 2}]
\addplot[fill=red!60, /pgf/bar shift=-6pt] table[x=image, y=part2_1] {\speedupVTwo};
\addplot[fill=blue!60, /pgf/bar shift=0pt] table[x=image, y=part2_2] {\speedupVTwo};
\addplot[fill=green!50!black!60, /pgf/bar shift=6pt] table[x=image, y=part2_3] {\speedupVTwo};
\addplot[only marks, mark=-, mark size=10pt, very thick, black, /pgf/bar shift=0pt]
  table[x=image, y=bound] {\speedupVTwo};
\end{groupplot}
\end{tikzpicture}\\[4pt]
\ref{legSpeedup}
\caption{Speedup over the serial version ($k=\cfgK$, $N=\cfgRing$). The dashed line marks
speedup~1 (as fast as the serial program); the black marks show the bound
$T_\text{serial}/T_\text{S1}$ for each image. Exact values are in Table~\ref{tab:speedup}.}
\label{fig:speedup}
\end{figure}

\paragraph{Variant 2 reaches the bound.} For all three images, every transport reaches a speedup
of $1.06$--$1.08\times$, within about 1\% of the bound (Figure~\ref{fig:speedup}). Three values (all
with threads) lie slightly \emph{above} the bound, which is within the run-to-run variation. For image~4
with threads, the pipeline takes 37.51\,s, while S1 alone needs 37.46\,s of computing in the same
run. Almost all of S2's and S3's work is hidden behind S1.

\paragraph{Variant 1 is slower than serial.} In Variant~1 all versions are 2--8\% \emph{slower}
than the serial program. The cause is verification on the bottleneck stage. In Variant~1, S1
computes a checksum for every strip of every iteration, which the serial program does not do.
For image~4 with threads, S1 spends 37.77\,s computing and 4.87\,s computing checksums
(Table~\ref{tab:stages4}). Together that is 42.64\,s, almost exactly the pipeline time of
42.66\,s. S2 and S3 also verify every strip, but they have spare time and that cost is hidden. In
Variant~2 only one pass in \cfgPasses{} is sent, so S1's checksum time drops to 0.05\,s. Without
checksums, Variant~1 would be expected to reach about the same speedup as Variant~2.

\paragraph{Comparing the transports.} Threads are the fastest in every row (tied with shared
memory for image~1, Variant~1). The differences are small (at most 3\%), because S1, not
communication, limits the pipeline. In Variant~1, shared memory comes second and pipes last:
every strip is copied into and out of the kernel, and S1 waits for the pipe (\emph{wait out}
0.88\,s for image~4, against 0.01\,s or less with the rings). In Variant~2, pipes and shared
memory are within 0.5\% of each other and swap places between images.

\subsection{Stage breakdown}
\label{sec:stages}
Each stage measures where its time goes (seconds):
\begin{itemize}
\item \emph{busy}: time in its kernel;
\item \emph{verify}: time computing checksums (sealing outgoing strips and checking incoming ones);
\item \emph{wait in}: time waiting for input (for pipes, including the copy);
\item \emph{wait out}: time waiting for space in the output channel (for pipes, including the copy);
\item \emph{first in}: time from the start until the first strip arrived (ms).
\end{itemize}
\begin{table}[htbp]
\centering
\caption{Stage breakdown for image 4 ($k=\cfgK$, $N=\cfgRing$). Images 1 and 2 are in
Appendix~\ref{app:stages}.}
\label{tab:stages4}
\begin{tabular}{llrrrrrr}
\toprule
Variant & Program & Stage & Busy & Verify & Wait in & Wait out & First in \\
\midrule
1 & Threads & S1 & 37.77 & 4.87 & 0.00 & 0.01 & -- \\
 &  & S2 & 1.31 & 9.74 & 31.57 & 0.01 & 1.34 \\
 &  & S3 & 1.61 & 4.87 & 36.16 & 0.00 & 1.47 \\
\addlinespace
1 & Pipes & S1 & 38.08 & 4.88 & 0.00 & 0.88 & -- \\
 &  & S2 & 1.56 & 11.80 & 29.53 & 0.90 & 1.07 \\
 &  & S3 & 2.30 & 7.62 & 33.82 & 0.00 & 1.42 \\
\addlinespace
1 & Shared mem. & S1 & 37.84 & 4.88 & 0.00 & 0.00 & -- \\
 &  & S2 & 1.91 & 14.45 & 25.82 & 0.04 & 1.10 \\
 &  & S3 & 2.17 & 7.28 & 33.37 & 0.00 & 1.31 \\
\addlinespace
2 & Threads & S1 & 37.46 & 0.05 & 0.00 & 0.00 & -- \\
 &  & S2 & 1.27 & 0.10 & 32.54 & 0.00 & 0.72 \\
 &  & S3 & 1.52 & 0.05 & 32.36 & 0.00 & 0.87 \\
\addlinespace
2 & Pipes & S1 & 37.80 & 0.05 & 0.00 & 0.01 & -- \\
 &  & S2 & 1.28 & 0.11 & 32.83 & 0.01 & 1.05 \\
 &  & S3 & 1.75 & 0.07 & 32.47 & 0.00 & 1.26 \\
\addlinespace
2 & Shared mem. & S1 & 37.89 & 0.05 & 0.00 & 0.00 & -- \\
 &  & S2 & 1.36 & 0.13 & 32.82 & 0.00 & 1.04 \\
 &  & S3 & 1.71 & 0.08 & 32.61 & 0.00 & 1.34 \\
\bottomrule
\end{tabular}
\end{table}
The breakdown confirms that S1 is the bottleneck. S1 hardly ever waits: \emph{wait in} is
always zero because it reads the image directly, and \emph{wait out} stays small (at most 0.88\,s,
with pipes) because the downstream stages keep up. S2 and S3 spend most of the run waiting for
input. The pipeline fills quickly: S3 receives its first strip after 0.9--1.5\,ms, out of a run
of about 40\,s. In Variant~1, S2 has the highest verify time, because it both checks every incoming strip
and seals every outgoing one.

\subsection{Effect of the strip height \texorpdfstring{$k$}{k}}
Pipeline time in seconds for $k=1,2,4,\dots,128$ and $N=\cfgRing$ (T = threads, P = pipes,
S = shared memory). The lowest time in each column is in bold.

\begin{table}[htbp]
\centering
\caption{Variant 1: pipeline time (seconds) against the strip height $k$.}
\label{tab:k-v1}
\begin{tabular}{r rrr rrr rrr}
\toprule
& \multicolumn{3}{c}{Image 1} & \multicolumn{3}{c}{Image 2} & \multicolumn{3}{c}{Image 4} \\
\cmidrule(lr){2-4}\cmidrule(lr){5-7}\cmidrule(lr){8-10}
$k$ & T & P & S & T & P & S & T & P & S \\
\midrule
1 & 3.05 & 3.27 & 3.31 & 11.34 & 11.93 & 11.81 & 42.77 & 43.70 & 43.95 \\
2 & 3.03 & 3.22 & 3.41 & 11.30 & \textbf{11.55} & 11.72 & \textbf{42.55} & 43.62 & 43.40 \\
4 & \textbf{2.97} & 3.07 & 3.10 & 11.32 & 11.60 & 11.45 & 42.60 & \textbf{43.28} & 43.17 \\
8 & 2.98 & \textbf{3.05} & 2.99 & 11.29 & 11.57 & 11.47 & 42.61 & 43.47 & 43.01 \\
16 & 2.99 & \textbf{3.05} & 3.08 & 11.28 & 11.62 & 11.40 & \textbf{42.55} & 43.77 & 42.87 \\
32 & 3.01 & 3.13 & 2.96 & \textbf{11.23} & 16.34 & 11.37 & 42.61 & 57.67 & \textbf{42.66} \\
64 & 2.99 & 4.32 & 2.98 & 11.25 & 15.39 & 11.39 & 42.69 & 60.43 & 42.77 \\
128 & 2.98 & 5.45 & \textbf{2.95} & 11.30 & 13.73 & \textbf{11.31} & 42.59 & 56.17 & 42.78 \\
\bottomrule
\end{tabular}
\end{table}

\begin{table}[htbp]
\centering
\caption{Variant 2: pipeline time (seconds) against the strip height $k$.}
\label{tab:k-v2}
\begin{tabular}{r rrr rrr rrr}
\toprule
& \multicolumn{3}{c}{Image 1} & \multicolumn{3}{c}{Image 2} & \multicolumn{3}{c}{Image 4} \\
\cmidrule(lr){2-4}\cmidrule(lr){5-7}\cmidrule(lr){8-10}
$k$ & T & P & S & T & P & S & T & P & S \\
\midrule
1 & 2.62 & 2.66 & 2.72 & 9.98 & 10.09 & 10.08 & 37.85 & \textbf{37.78} & 37.78 \\
2 & 2.67 & 2.65 & 2.65 & 9.96 & 10.10 & 10.05 & \textbf{37.59} & 37.97 & \textbf{37.72} \\
4 & 2.64 & 2.65 & 2.66 & 10.01 & \textbf{10.03} & \textbf{9.99} & 37.70 & 37.84 & 37.92 \\
8 & 2.63 & \textbf{2.64} & 2.70 & \textbf{9.93} & \textbf{10.03} & \textbf{9.99} & 37.66 & 37.94 & 37.80 \\
16 & \textbf{2.61} & 2.65 & 2.68 & 9.96 & 10.04 & 10.07 & 37.64 & 37.96 & 37.91 \\
32 & 2.64 & \textbf{2.64} & \textbf{2.63} & 10.00 & 10.11 & 10.05 & 37.63 & 37.90 & 37.75 \\
64 & 2.65 & 2.76 & 2.68 & 10.01 & \textbf{10.03} & 10.03 & 37.77 & 37.84 & 37.74 \\
128 & 2.64 & 2.71 & 2.65 & 9.95 & 10.06 & 10.02 & 37.75 & 37.87 & 37.83 \\
\bottomrule
\end{tabular}
\end{table}

\begin{figure}[p]
\centering
\begin{tikzpicture}
\begin{groupplot}[
  group style={group size=2 by 3, horizontal sep=1.6cm, vertical sep=1.4cm,
               xlabels at=edge bottom, ylabels at=edge left},
  sweep plot, cycle list name=programs,
  xtick={1,2,4,8,16,32,64,128}, xlabel={$k$ (rows per strip)},
  ylabel={Relative to threads, $k=\cfgK$}, ymin=0.9, ymax=1.9,
  extra y ticks={1}, extra y tick labels={},
  extra y tick style={grid=major, major grid style={black, dashed}},
  legend columns=3, legend style={font=\small, /tikz/every even column/.append style={column sep=6pt}}
]
% All times in a panel divided by the threads time at k = 8 (values from the k tables).
\nextgroupplot[title={Image 1, Variant 1}, legend to name=legSweepB]
\addplot table[x=k, y expr=\thisrow{part2_1}/2.98] {\sweepBimgOneVOne};
\addplot table[x=k, y expr=\thisrow{part2_2}/2.98] {\sweepBimgOneVOne};
\addplot table[x=k, y expr=\thisrow{part2_3}/2.98] {\sweepBimgOneVOne};
\legend{Threads, Pipes, Shared mem.}
\nextgroupplot[title={Image 1, Variant 2}]
\addplot table[x=k, y expr=\thisrow{part2_1}/2.63] {\sweepBimgOneVTwo};
\addplot table[x=k, y expr=\thisrow{part2_2}/2.63] {\sweepBimgOneVTwo};
\addplot table[x=k, y expr=\thisrow{part2_3}/2.63] {\sweepBimgOneVTwo};
\nextgroupplot[title={Image 2, Variant 1}]
\addplot table[x=k, y expr=\thisrow{part2_1}/11.29] {\sweepBimgTwoVOne};
\addplot table[x=k, y expr=\thisrow{part2_2}/11.29] {\sweepBimgTwoVOne};
\addplot table[x=k, y expr=\thisrow{part2_3}/11.29] {\sweepBimgTwoVOne};
\nextgroupplot[title={Image 2, Variant 2}]
\addplot table[x=k, y expr=\thisrow{part2_1}/9.93] {\sweepBimgTwoVTwo};
\addplot table[x=k, y expr=\thisrow{part2_2}/9.93] {\sweepBimgTwoVTwo};
\addplot table[x=k, y expr=\thisrow{part2_3}/9.93] {\sweepBimgTwoVTwo};
\nextgroupplot[title={Image 4, Variant 1}]
\addplot table[x=k, y expr=\thisrow{part2_1}/42.61] {\sweepBimgFourVOne};
\addplot table[x=k, y expr=\thisrow{part2_2}/42.61] {\sweepBimgFourVOne};
\addplot table[x=k, y expr=\thisrow{part2_3}/42.61] {\sweepBimgFourVOne};
\nextgroupplot[title={Image 4, Variant 2}]
\addplot table[x=k, y expr=\thisrow{part2_1}/37.66] {\sweepBimgFourVTwo};
\addplot table[x=k, y expr=\thisrow{part2_2}/37.66] {\sweepBimgFourVTwo};
\addplot table[x=k, y expr=\thisrow{part2_3}/37.66] {\sweepBimgFourVTwo};
\end{groupplot}
\end{tikzpicture}\\[6pt]
\ref{legSweepB}
\caption{Effect of the strip height $k$ ($N=\cfgRing$). In each panel, all pipeline times are
divided by the time of the threads version at the default $k=\cfgK$ (dashed line), so the shape of
a line shows the effect of $k$ and the distance between lines shows the difference between the
programs. All panels use the same scale. In Variant~2 (right) only \cfgResults{} images pass
through the pipeline, so $k$ has no visible effect and all lines stay at about 1. Absolute times
are in Tables~\ref{tab:k-v1} and~\ref{tab:k-v2}.}
\label{fig:k}
\end{figure}

\paragraph{Threads and shared memory.} For the threads version, the time for images 2 and 4
changes by less than 1\% across all values of $k$, and for the short runs of image~1 by up to
2.7\% (2.97--3.05\,s in Variant~1), without a trend in $k$. This is within the run-to-run
variation. Shared memory is a few percent slower at the smallest strips: image~4 takes
43.95\,s at $k=1$ against 42.66\,s at $k=32$, because every strip costs two semaphore operations
on each side, and $k=1$ has 32 times as many strips as $k=32$. Above $k\approx16$ shared memory
matches threads.

\paragraph{Pipes and the pipe buffer.} In Variant~1, the pipes version slows down sharply
compared with $k=8$: by up to 39\% for image~4 (43.5\,s at $k=8$, 60.4\,s at $k=64$), 41\% for
image~2 and 79\% for image~1 (3.05\,s at $k=8$, 5.45\,s at $k=128$). It does so exactly when a strip
becomes larger than the default Linux pipe capacity of 65{,}536 bytes:
\begin{center}
\begin{tabular}{rrrrr}
\toprule
Image & Bytes per row & Largest tested $k$ that fits & Smallest that does not & Time (s) \\
\midrule
1 & 1{,}350 & 32 (43{,}200 B) & 64 (86{,}400 B) & 3.13 $\to$ 4.32 \\
2 & 2{,}358 & 16 (37{,}728 B) & 32 (75{,}456 B) & 11.62 $\to$ 16.34 \\
4 & 3{,}840 & 16 (61{,}440 B) & 32 (122{,}880 B) & 43.77 $\to$ 57.67 \\
\bottomrule
\end{tabular}
\end{center}
While a strip fits into the pipe, the writer copies it in one go and continues computing. A
larger strip can only be written once the reader has started draining the pipe, so producer and
consumer have to take turns within every strip. For image~2, S1's \emph{wait out} rises from
0.24\,s at $k=16$ to 2.32\,s at $k=32$. S1's \emph{busy} time also rises, from 10.08\,s to
12.49\,s, although it does the same computation. A plausible cause is that S1 now blocks and is
woken up thousands of times, which costs it cache contents and possibly a fast CPU core; we did
not investigate this further. Variant~2 sends only \cfgResults{} images, so this effect is not
visible there.

\paragraph{Variant 2.} In Variant~2 the strip height makes no systematic difference for any
transport. For images 2 and 4 all values lie within 1\%; the short runs of image~1 vary more
(up to 4.5\%, pipes: 2.64--2.76\,s) but without a trend in $k$. The best values in
Table~\ref{tab:k-v2} are therefore mostly noise.

\subsection{Effect of the ring capacity \texorpdfstring{$N$}{N}}
Pipeline time in seconds for $N=1,2,4,\dots,64$ and $k=\cfgK$ (T = threads, S = shared memory;
pipes have no ring). The lowest time in each column is in bold.

\begin{table}[htbp]
\centering
\caption{Variant 1: pipeline time (seconds) against the ring capacity $N$.}
\label{tab:N-v1}
\begin{tabular}{r rr rr rr}
\toprule
& \multicolumn{2}{c}{Image 1} & \multicolumn{2}{c}{Image 2} & \multicolumn{2}{c}{Image 4} \\
\cmidrule(lr){2-3}\cmidrule(lr){4-5}\cmidrule(lr){6-7}
$N$ & T & S & T & S & T & S \\
\midrule
1 & 3.78 & 12.00 & 14.30 & 42.95 & 54.39 & 129.86 \\
2 & 2.99 & 3.09 & 11.39 & 13.51 & 42.36 & 43.35 \\
4 & \textbf{2.94} & 3.04 & \textbf{11.19} & 11.46 & 42.37 & \textbf{42.81} \\
8 & 2.97 & 3.04 & 11.29 & \textbf{11.36} & \textbf{42.33} & 43.02 \\
16 & 2.98 & \textbf{3.00} & 11.24 & 11.39 & 42.69 & 43.00 \\
32 & 2.98 & 3.02 & 11.28 & 11.41 & 42.61 & 43.06 \\
64 & 2.98 & \textbf{3.00} & 11.28 & 11.45 & 42.50 & 42.95 \\
\bottomrule
\end{tabular}
\end{table}

\begin{table}[htbp]
\centering
\caption{Variant 2: pipeline time (seconds) against the ring capacity $N$.}
\label{tab:N-v2}
\begin{tabular}{r rr rr rr}
\toprule
& \multicolumn{2}{c}{Image 1} & \multicolumn{2}{c}{Image 2} & \multicolumn{2}{c}{Image 4} \\
\cmidrule(lr){2-3}\cmidrule(lr){4-5}\cmidrule(lr){6-7}
$N$ & T & S & T & S & T & S \\
\midrule
1 & 2.65 & 3.01 & 10.02 & 10.28 & 37.79 & 38.92 \\
2 & 2.64 & 2.71 & 10.05 & 10.03 & \textbf{37.61} & 37.93 \\
4 & 2.65 & 2.69 & 10.00 & 10.05 & 37.93 & 37.89 \\
8 & \textbf{2.62} & \textbf{2.64} & 10.03 & \textbf{10.00} & 37.62 & 38.06 \\
16 & 2.66 & 2.70 & 10.00 & 10.01 & 37.95 & 37.99 \\
32 & \textbf{2.62} & 2.65 & 10.01 & 10.09 & 37.83 & 37.79 \\
64 & 2.66 & 2.65 & \textbf{9.98} & 10.01 & 37.74 & \textbf{37.73} \\
\bottomrule
\end{tabular}
\end{table}

\begin{figure}[p]
\centering
\begin{tikzpicture}
\begin{groupplot}[
  group style={group size=2 by 3, horizontal sep=1.6cm, vertical sep=1.4cm,
               xlabels at=edge bottom, ylabels at=edge left},
  sweep plot, cycle list name=rings,
  xtick={1,2,4,8,16,32,64}, xlabel={$N$ (ring capacity, strips)},
  ylabel={Relative to threads, $N=\cfgRing$}, ymin=0.8, ymax=4.2, ytick={1,2,3,4},
  extra y ticks={1}, extra y tick labels={},
  extra y tick style={grid=major, major grid style={black, dashed}},
% All times in a panel divided by the threads time at N = 16 (values from the N tables).
]
% Each program's time divided by its own time at N = 16 (values from the N tables).
\nextgroupplot[title={Image 1, Variant 1}, legend to name=legSweepC]
\addplot table[x=N, y expr=\thisrow{part2_1}/2.98] {\sweepCimgOneVOne};
\addplot table[x=N, y expr=\thisrow{part2_3}/2.98] {\sweepCimgOneVOne};
\legend{Threads, Shared mem.}
\nextgroupplot[title={Image 1, Variant 2}]
\addplot table[x=N, y expr=\thisrow{part2_1}/2.66] {\sweepCimgOneVTwo};
\addplot table[x=N, y expr=\thisrow{part2_3}/2.66] {\sweepCimgOneVTwo};
\nextgroupplot[title={Image 2, Variant 1}]
\addplot table[x=N, y expr=\thisrow{part2_1}/11.24] {\sweepCimgTwoVOne};
\addplot table[x=N, y expr=\thisrow{part2_3}/11.24] {\sweepCimgTwoVOne};
\nextgroupplot[title={Image 2, Variant 2}]
\addplot table[x=N, y expr=\thisrow{part2_1}/10.00] {\sweepCimgTwoVTwo};
\addplot table[x=N, y expr=\thisrow{part2_3}/10.00] {\sweepCimgTwoVTwo};
\nextgroupplot[title={Image 4, Variant 1}]
\addplot table[x=N, y expr=\thisrow{part2_1}/42.69] {\sweepCimgFourVOne};
\addplot table[x=N, y expr=\thisrow{part2_3}/42.69] {\sweepCimgFourVOne};
\nextgroupplot[title={Image 4, Variant 2}]
\addplot table[x=N, y expr=\thisrow{part2_1}/37.95] {\sweepCimgFourVTwo};
\addplot table[x=N, y expr=\thisrow{part2_3}/37.95] {\sweepCimgFourVTwo};
\end{groupplot}
\end{tikzpicture}\\[6pt]
\ref{legSweepC}
\caption{Effect of the ring capacity $N$ ($k=\cfgK$). In each panel, all pipeline times are
divided by the time of the threads version at the default $N=\cfgRing$ (dashed line), so the shape
of a line shows the effect of $N$ and the distance between lines shows the difference between
the programs. All panels use the same scale. In Variant~2 (right) only \cfgResults{} images pass
through the pipeline, so $N$ has almost no effect (at most 11\% at $N=1$). Absolute times are in
Tables~\ref{tab:N-v1} and~\ref{tab:N-v2}.}
\label{fig:N}
\end{figure}

\paragraph{$N \ge 2$.} For threads, the ring capacity has no measurable effect from two slots
upwards, and for shared memory from four slots upwards. S1 is the bottleneck, so the rings are
almost always close to empty and extra slots are rarely used. The one exception is shared memory
with two slots in Variant~1: image~2 takes 13.51\,s against 11.36--11.46\,s for $N\ge4$ (image~1:
3.09 against 3.00--3.04\,s).

\paragraph{$N = 1$.} With a single slot, S1 cannot start the next strip until S2 has released the
previous one, so the stages run in lockstep instead of overlapping. Threads lose up to 28\%
(image~4: 54.4\,s against 42.4\,s at $N=2$). Shared memory is three times slower (129.9\,s
against 43.4\,s). With shared memory, S1's busy time for image~2 more than doubles, from 11.6\,s at
$N=2$ to 22.2\,s at $N=1$, even though the computation is the same. The difference between the
two transports is in how they wait: a thread spins and yields and stays on its core, while a
process blocked on a semaphore sleeps and must be woken by the kernel for every strip. We suspect
these thousands of sleep/wake-up cycles cost S1 its cache contents and CPU frequency, but we did
not verify this. In Variant~2 the effect is much smaller (at most 3\% for images 2 and 4, 11\%
for the small image~1), because only \cfgResults{} images pass through the ring.

% =====================================================================
\section{Results: Part III (two machines)}
Machine~A is a laptop running Linux under WSL2 on Windows; machine~B is a laptop running Linux
natively. Both were connected to the same Wi-Fi hotspot. All runs use $k=\cfgK$ and $N=\cfgRing$
and the same three images, and all times are end-to-end times measured on A.

Both partitions send the same amount of data over the network: one image-sized stream of strips
per iteration (Variant~1) or per result (Variant~2). They differ only in where the stages run:
\begin{itemize}
\item \textbf{Partition 1:} A runs S1 and S2 (two threads with a ring between them), B runs S3.
  The S2$\to$S3 link crosses the network.
\item \textbf{Partition 2:} A runs S1, B runs S2 and S3. The S1$\to$S2 link crosses the network.
\end{itemize}

\paragraph{Correctness.} In all 30 runs, both ends of the network link reported identical strip
counts, byte counts and stream hashes, so every strip arrived intact and in order. For image~4 the
stream hashes are also identical to those of the Part~II runs on a single machine
(Table~\ref{tab:links}): the same strips were sent, in the same order.

\begin{table}[htbp]
\centering
\caption{End-to-end time on machine A (seconds) for both partitions, and speedup over the serial
program run on machine A. Variant~1: one run each. Variant~2 and serial: fastest of three or more
runs, all in the same session. Best per row in bold.}
\label{tab:part3}
\begin{tabular}{rr r rr rr}
\toprule
& & & \multicolumn{2}{c}{Partition 1} & \multicolumn{2}{c}{Partition 2} \\
\cmidrule(lr){4-5}\cmidrule(lr){6-7}
Image & Variant & Serial (A) & Time & Speedup & Time & Speedup \\
\midrule
1 & 1 & 2.48 & \textbf{47.52} & \textbf{0.052} & 49.71 & 0.050 \\
1 & 2 & 2.48 & \textbf{2.31} & \textbf{1.077} & 2.34 & 1.062 \\
\addlinespace
2 & 1 & 9.42 & 247.62 & 0.038 & \textbf{205.38} & \textbf{0.046} \\
2 & 2 & 9.42 & \textbf{9.08} & \textbf{1.038} & 9.23 & 1.021 \\
\addlinespace
4 & 1 & 36.29 & 885.02 & 0.041 & \textbf{750.12} & \textbf{0.048} \\
4 & 2 & 36.29 & \textbf{37.78} & \textbf{0.961} & 37.87 & 0.958 \\
\bottomrule
\end{tabular}
\end{table}

\paragraph{Measuring speedup on machine A.} Speedups are measured against the serial program on
machine A, because A runs S1, the stage that sets the pace. A runs Linux under Windows (WSL2),
where Windows decides the clock speed and which core a program runs on, and A's speed turned out
to differ between sessions: in one earlier session, S1 took 4.1\,s for image~1, against 2.3\,s in
the session used here. Comparing a serial run from one session with a pipeline run from another
would therefore give meaningless speedups. For Variant~2, the serial program and both partitions
were therefore run alternately in the same session, three rounds each, and the fastest run of
each is used. Interference can only make a run slower, so the fastest run is the best estimate of
the machine's actual speed. Within this session S1's compute time varied far less than between
sessions (2.25--2.43\,s for image~1, 32.8--35.3\,s for image~4, over serial and pipeline runs),
although for image~4 it was somewhat higher in the pipeline runs (33.4--35.3\,s) than in the
serial runs (32.8--34.0\,s). Variant~1 runs the same
computation as Variant~2 (on machine~B the serial times of the two variants differ by less than
3\%), so the same serial times are used for it. Its speedups are so far below 1 that the
difference does not matter.

\paragraph{Speedup.} On machine A, S1 takes about 90\% of the serial time, so the bound
$T_\text{serial}/T_\text{S1}$ is about $1.11\times$ (image~4: $36.29/32.78 = 1.107$). In
Variant~1 every configuration is 19 to 26 times \emph{slower} than serial (speedups
0.038--0.052), because of the network (see below). In Variant~2, images 1 and 2 gain from the
pipeline (up to $1.077\times$ for image~1, close to its bound of $1.106\times$), but image~4 is
about 4\% slower than serial.

\paragraph{Variant 1 is limited by the network.} In Variant~1 every iteration sends a whole image
across the network, 3.5\,GB in total for image~4. Dividing the bytes sent by the run time gives the
rate the link sustained (Table~\ref{tab:part3-rate}): 3.7--5.1\,MB/s, or about 30--40\,Mbit/s,
which is typical for a shared Wi-Fi hotspot. The stage statistics show where the time goes. For
image~4 in partition~1, S1 computed for 34.8\,s and spent 846\,s, 96\% of the run, waiting for
space in its output channel: the pipeline ran at the speed of the network, not of S1. The same
run on one machine took 42.7\,s (Table~\ref{tab:speedup}), about 20 times faster. In
partition~2 the picture is the same: S1 waited for 707\,s of the 750\,s run (94\%). Since S1
spends almost the whole run waiting in both partitions, the difference between them (885 against
750\,s for image~4) most likely reflects how fast the shared hotspot carried data during each run,
rather than where S2 runs.

\begin{table}[htbp]
\centering
\caption{Variant 1: data sent over the network (headers and pixels) and the rate achieved
(bytes divided by the end-to-end time; MB = $10^6$ bytes).}
\label{tab:part3-rate}
\begin{tabular}{rr rr rr}
\toprule
& & \multicolumn{2}{c}{Partition 1} & \multicolumn{2}{c}{Partition 2} \\
\cmidrule(lr){3-4}\cmidrule(lr){5-6}
Image & Data (MB) & Time (s) & Rate (MB/s) & Time (s) & Rate (MB/s) \\
\midrule
1 & 243.6 & 47.52 & 5.13 & 49.71 & 4.90 \\
2 & 927.9 & 247.62 & 3.75 & 205.38 & 4.52 \\
4 & 3{,}497.1 & 885.02 & 3.95 & 750.12 & 4.66 \\
\bottomrule
\end{tabular}
\end{table}

\paragraph{Variant 2 hides the network for small results.} Variant~2 sends only \cfgResults{}
results, and between two results S1 computes \cfgPasses{} passes locally, so there is plenty of
time to transfer a result in the background. What matters is whether S1 can hand over a whole
result without waiting: the strips of pass~0 are produced much faster than the network can carry
them, and they first have to fit into the buffers between S1 and the network (the socket buffers,
and in partition~1 also the ring between S1 and S2). One result is 0.24\,MB for image~1 and
0.93\,MB for image~2. These fit, and S1 waited at most 0.5\,s in total over a run. For image~4 a
result is 3.5\,MB. It does not fit, so S1 is blocked on every result until enough has been sent;
this added 2.6--5.7\,s of waiting per run. The pipeline can save at most 3.5\,s
($T_\text{serial}-T_\text{S1}$), and in addition S1 computed somewhat more slowly in the pipeline
runs than in the serial runs (see above). In the fastest partition~1 run, for example, S1
computed for 34.95\,s and waited 2.78\,s, for 37.78\,s in total against 36.29\,s serial. Both partitions behave the same: over three rounds, the fastest
times differ by less than 2\% for every image, and S1's waiting varied as much within each
partition as between them.

% =====================================================================
\section{Conclusion}
All versions produce output identical to the serial reference, and in every run both ends of
every link agreed on the number, content and order of the strips. The achievable speedup of
this three-stage pipeline is limited to about $1.08\times$, because S1 accounts for about 93\% of
the work. In Variant~2 all three transports reach this bound to within about 1\%. In Variant~1
the cost of verifying every strip falls on S1, the bottleneck stage, and makes the pipelined
versions 2--8\% slower than the serial program. Threads with a lock-free ring are the fastest
transport; in Variant~1 shared memory comes second and pipes last. Two parameters matter only at
their extremes: pipes slow down by 39--79\% once a strip no longer fits into the 64\,KiB pipe
buffer, and a ring of a single slot forces the stages into lockstep, which costs a factor of three
to four with semaphores. A pipeline that splits S1 itself across several workers (for example
by rows) would be needed for a larger speedup.

Across two machines, the data again arrived exactly as sent in every run. Variant~1 sends every
iteration over the network and is limited by it: over a Wi-Fi hotspot at about 4--5\,MB/s, the
largest image took 750--885\,s, about 20 times longer than on one machine. Variant~2 sends
100 times less data. For the two smaller images it still gains up to $1.08\times$ over the serial
program, but for the largest image one result no longer fits into the buffers on the way, S1 has
to wait for the network, and the pipeline ends up about 4\% slower than serial. Where the network
link lies (partition 1 or 2) made no measurable difference.

% =====================================================================
\clearpage
\appendix
\section{Stage breakdown for images 1 and 2}
\label{app:stages}

\begin{table}[htbp]
\centering
\caption{Stage breakdown for image 1 ($k=\cfgK$, $N=\cfgRing$), seconds; first in in ms.}
\label{tab:stages1}
\begin{tabular}{llrrrrrr}
\toprule
Variant & Program & Stage & Busy & Verify & Wait in & Wait out & First in \\
\midrule
1 & Threads & S1 & 2.64 & 0.34 & 0.00 & 0.00 & -- \\
 &  & S2 & 0.09 & 0.69 & 2.21 & 0.00 & 0.60 \\
 &  & S3 & 0.11 & 0.34 & 2.53 & 0.00 & 0.65 \\
\addlinespace
1 & Pipes & S1 & 2.65 & 0.34 & 0.00 & 0.06 & -- \\
 &  & S2 & 0.10 & 0.72 & 2.16 & 0.07 & 0.72 \\
 &  & S3 & 0.16 & 0.50 & 2.39 & 0.00 & 0.92 \\
\addlinespace
1 & Shared mem. & S1 & 2.62 & 0.34 & 0.00 & 0.00 & -- \\
 &  & S2 & 0.12 & 0.90 & 1.91 & 0.00 & 0.68 \\
 &  & S3 & 0.15 & 0.49 & 2.34 & 0.00 & 0.84 \\
\addlinespace
2 & Threads & S1 & 2.65 & 0.00 & 0.00 & 0.00 & -- \\
 &  & S2 & 0.09 & 0.01 & 2.31 & 0.00 & 0.45 \\
 &  & S3 & 0.11 & 0.00 & 2.29 & 0.00 & 0.50 \\
\addlinespace
2 & Pipes & S1 & 2.67 & 0.00 & 0.00 & 0.00 & -- \\
 &  & S2 & 0.09 & 0.01 & 2.32 & 0.00 & 0.60 \\
 &  & S3 & 0.12 & 0.00 & 2.31 & 0.00 & 0.75 \\
\addlinespace
2 & Shared mem. & S1 & 2.68 & 0.00 & 0.00 & 0.00 & -- \\
 &  & S2 & 0.09 & 0.01 & 2.33 & 0.00 & 0.68 \\
 &  & S3 & 0.15 & 0.01 & 2.27 & 0.00 & 0.79 \\
\bottomrule
\end{tabular}
\end{table}

\begin{table}[htbp]
\centering
\caption{Stage breakdown for image 2 ($k=\cfgK$, $N=\cfgRing$), seconds; first in in ms.}
\label{tab:stages2}
\begin{tabular}{llrrrrrr}
\toprule
Variant & Program & Stage & Busy & Verify & Wait in & Wait out & First in \\
\midrule
1 & Threads & S1 & 9.98 & 1.29 & 0.00 & 0.00 & -- \\
 &  & S2 & 0.34 & 2.61 & 8.32 & 0.00 & 0.59 \\
 &  & S3 & 0.42 & 1.30 & 9.55 & 0.00 & 0.72 \\
\addlinespace
1 & Pipes & S1 & 10.03 & 1.31 & 0.00 & 0.28 & -- \\
 &  & S2 & 0.42 & 3.19 & 7.73 & 0.26 & 1.01 \\
 &  & S3 & 0.58 & 1.86 & 9.16 & 0.00 & 1.24 \\
\addlinespace
1 & Shared mem. & S1 & 10.05 & 1.30 & 0.00 & 0.00 & -- \\
 &  & S2 & 0.42 & 3.18 & 7.66 & 0.01 & 1.20 \\
 &  & S3 & 0.62 & 2.06 & 8.71 & 0.00 & 1.54 \\
\addlinespace
2 & Threads & S1 & 9.96 & 0.01 & 0.00 & 0.00 & -- \\
 &  & S2 & 0.36 & 0.03 & 8.64 & 0.00 & 0.78 \\
 &  & S3 & 0.43 & 0.01 & 8.59 & 0.00 & 0.86 \\
\addlinespace
2 & Pipes & S1 & 10.04 & 0.01 & 0.00 & 0.00 & -- \\
 &  & S2 & 0.34 & 0.03 & 8.72 & 0.00 & 0.67 \\
 &  & S3 & 0.48 & 0.02 & 8.63 & 0.00 & 0.92 \\
\addlinespace
2 & Shared mem. & S1 & 10.00 & 0.01 & 0.00 & 0.00 & -- \\
 &  & S2 & 0.36 & 0.03 & 8.67 & 0.00 & 0.90 \\
 &  & S3 & 0.42 & 0.01 & 8.63 & 0.00 & 1.09 \\
\bottomrule
\end{tabular}
\end{table}

\end{document}
